% Esquisse d'un programme (cote 119, copy 2): archive cote 119, site page 92 (PDF page 93), typescript page 14.
% typescript. Draft compared: esquisse/pages/esq_119_016.tex (copy 1, site 16).
% AI-transcribed by Claude Opus 5.5 (claude-opus-5-5, adaptive thinking, effort high), Message Batches API, batch esq119_full (msgbatch_015xqDHkTumDysRF1qngZhRJ), 2026-11-13; not yet checked by eye.
\scanpage{119/92}
\mspage{14}

s'est engagée dans une directions plus conceptuelle, pour arriver
à appréhender la nature de cette action de $\mathrm{I\!\Gamma}$. On s'aperçoit
d'emblée que grosso-modo cette action est exprimée par une cer-%
taine action ``extérieure'' de $\mathrm{I\!\Gamma}$ sur le compactifié profini du
groupe cartographique orienté $\underline{C}_2^{+}$ , et cette action à son tour
est déduite par passage au quotient de l'action extérieure cano-%
nique de $\mathrm{I\!\Gamma}$ sur le groupe fondamental profini $\hat{\pi}_{0,3}$ de $(U_{0,3})_{\bar{\mathbb{Q}}}$ ,
où $U_{0,3}$ désigne la courbe-type de genre $0$ sur le corps premier
$\mathbb{Q}$, privée de trois points. C'est ainsi que mon attention s'est
portée vers ce que j'ai appelé depuis la ``\uline{géométrie algébrique}
\uline{anabélienne}'', dont le point de départ est justement une étude
(pour le moment limitée à la caractéristique zéro) de l'action
de groupes de Galois ``absolus'' (notamment les groupes $\mathrm{Gal}(\bar{K}/K)$,
où $K$ est une extension de type fini du corps premier) sur des
groupes fondamentaux géométriques (profinis) de variétés algé-%
briques (définies sur $K$), et plus particulièrement (rompant avec
une tradition bien enracinée) des groupes fondamentaux qui sont
très éloignés des groupes abéliens (et que pour cette raison je
nomme ``\uline{anabéliens}''). Parmi ces groupes, et très proche du groupe
$\hat{\pi}_{0,3}$ , il y a le compactifié profini du groupe modulaire $S\ell(2,\mathbb{Z})$,
dont le quotient par le centre $\pm 1$ contient le précédent comme
sous-groupe de congruence mod 2, et peut s'interpréter d'ailleurs
également comme un groupe ``cartographique'' orienté, savoir celui
qui classifie les cartes orientées \uline{triangulées} (i.e. celles dont
les faces sont des triangles ou des monogones).

Toute carte finie orientée donne lieu à une courbe al-%
gébrique projective et lisse définie sur $\bar{\mathbb{Q}}$ , et il se pose alors
immédiatement la question : quelles sont les courbes algébriques
sur $\bar{\mathbb{Q}}$ obtenues ainsi - les obtiendrait-on toutes, qui sait ?
En termes plus savants, serait-il vrai que toute courbe algébrique
projective et lisse définie sur un corps de nombres interviendrait
comme une ``courbe modulaire'' possible pour paramétriser les cour-%
bes elliptiques munies d'une rigidification convenable ? Une telle
supposition avait l'air à tel point dingue que j'étais presque

\margin{92}
